\chapter{Aplicaciones lineales}\label{ch:b1:linmaps}

Las aplicaciones que merece la pena estudiar entre \omterm{def:b1:vspaces:def}{espacios
vectoriales} son las compatibles con la estructura: las
\emph{\omterm{def:b1:linmaps:def}{aplicaciones lineales}}. Sus dos \omterm{def:b1:vspaces:subspace}{subespacios} fundamentales
---núcleo e imagen--- miden la \omterm{def:b1:logic:inj}{inyectividad} y la \omterm{def:b1:logic:inj}{sobreyectividad}
y, en \omterm{def:b1:findim:def}{dimensión finita}, el teorema del \omterm{def:b1:linmaps:rank}{rango} ata sus tamaños en una
sola ley de conservación. Las proyecciones y las \omterm{def:b1:linmaps:projection}{simetrías}, y después
las \omterm{def:b1:linmaps:forms}{formas lineales} y los \omterm{def:b1:linmaps:forms}{hiperplanos}, cierran el capítulo.

\section{Definiciones y primeras propiedades}

\begin{definition}\label{def:b1:linmaps:def}
Sean $E, F$ \omterm{def:b1:vspaces:def}{espacios vectoriales} sobre $K$. Una \omterm{def:b1:logic:map}{aplicación}
$u \colon E \to F$ es \emph{lineal}\index{aplicación lineal} cuando
\[
\forall x, y \in E,\ \forall \lambda \in K, \qquad
u(x + \lambda y) = u(x) + \lambda u(y).
\]
Entonces $u(0) = 0$ y
$u(\sum \lambda_i x_i) = \sum \lambda_i u(x_i)$. El \omterm{def:b1:logic:sets}{conjunto}
$\mathcal{L}(E, F)$ de las aplicaciones lineales es él mismo un
\omterm{def:b1:vspaces:def}{espacio vectorial}; la composición de aplicaciones lineales es
lineal, y bilineal en la pareja. Un \emph{endomorfismo} es una
\omterm{def:b1:logic:map}{aplicación} lineal $E \to E$; un \emph{isomorfismo} es una
\omterm{def:b1:logic:map}{aplicación} lineal \omterm{def:b1:logic:inj}{biyectiva} (cuya inversa es entonces
automáticamente lineal); y $u^{-1}$ de un isomorfismo, así como las
composiciones de isomorfismos, son isomorfismos.
\end{definition}

\begin{proof}[Demostración de que la inversa es lineal]
Sea $u$ \omterm{def:b1:linmaps:def}{lineal} y \omterm{def:b1:logic:inj}{biyectiva}, sean $y, y' \in F$ y $\lambda \in K$.
Póngase $x = u^{-1}(y)$ y $x' = u^{-1}(y')$. Entonces
\[
u\bigl(x + \lambda x'\bigr) = u(x) + \lambda u(x')
= y + \lambda y' ,
\]
y, aplicando $u^{-1}$ a los dos extremos:
$u^{-1}(y + \lambda y') = x + \lambda x' = u^{-1}(y) +
\lambda\,u^{-1}(y')$. No se ha calculado nada de $u^{-1}$: la
linealidad se transporta solo por la propiedad que define $u$ --- un
patrón que conviene recordar, pues la estructura suele viajar gratis
a lomos de las biyecciones.
\end{proof}

\begin{proposition}[Una aplicación lineal se conoce en una base]\label{prop:b1:linmaps:basis}
Sea $(e_1, \dots, e_n)$ una \omterm{def:b1:vspaces:free}{base} de $E$ y sean $(v_1, \dots, v_n)$
vectores arbitrarios de $F$. Hay exactamente una \omterm{def:b1:linmaps:def}{aplicación lineal}
$u \colon E \to F$ con $u(e_i) = v_i$ para todo $i$. Además:
\[
u \text{ inyectiva} \iff (v_i) \text{ libre};
\qquad
u \text{ sobreyectiva} \iff (v_i) \text{ genera } F .
\]
\end{proposition}

\begin{proof}
Existencia y unicidad: todo $x$ tiene \omterm{prop:b1:vspaces:coordinates}{coordenadas} únicas
$x = \sum \lambda_i e_i$ (\cref{prop:b1:vspaces:coordinates}); la
linealidad fuerza $u(x) = \sum \lambda_i v_i$, y esa fórmula sí
define una \omterm{def:b1:linmaps:def}{aplicación lineal}.

\omterm{def:b1:logic:inj}{Inyectividad}: por la~\cref{prop:b1:linmaps:kernel} de más abajo, $u$
es \omterm{def:b1:logic:inj}{inyectiva} si y solo si su \omterm{def:b1:linmaps:kerim}{núcleo} es trivial. Ahora bien,
$u\bigl(\sum\lambda_i e_i\bigr) = 0$ significa exactamente
$\sum\lambda_i v_i = 0$. Si $(v_i)$ es \omterm{def:b1:vspaces:free}{libre}, esto fuerza que todos
los $\lambda_i = 0$, es decir, que el \omterm{def:b1:linmaps:kerim}{núcleo} se reduzca a $0$:
\omterm{def:b1:logic:inj}{inyectiva}. Y si $(v_i)$ es ligada, una relación no trivial
$\sum\lambda_i v_i = 0$ produce el vector no nulo
$\sum\lambda_i e_i$ en el \omterm{def:b1:linmaps:kerim}{núcleo} ($(e_i)$ es \omterm{def:b1:vspaces:free}{libre}): no \omterm{def:b1:logic:inj}{inyectiva}.
Las dos condiciones se corresponden término a término.

\omterm{def:b1:logic:inj}{Sobreyectividad}: la \omterm{def:b1:linmaps:kerim}{imagen} de $u$ es el \omterm{def:b1:logic:sets}{conjunto} de todos los
$\sum\lambda_i v_i$, es decir, exactamente
$\operatorname{Vect}(v_1, \dots, v_n)$, que es igual a $F$ si y solo
si la familia genera.
\end{proof}

\begin{definition}[Núcleo e imagen]\label{def:b1:linmaps:kerim}
Para $u \in \mathcal{L}(E, F)$:
\[
\ker u = \{x \in E : u(x) = 0\} \subseteq E,
\qquad
\operatorname{im} u = u(E) \subseteq F ,
\]
los dos \omterm{def:b1:vspaces:subspace}{subespacios} (comprobación directa con el
criterio).\index{núcleo}\index{imagen}
\end{definition}

\begin{method}[Núcleo e imagen, en la práctica]
\label{met:b1:linmaps:kerim}
\emph{\omterm{def:b1:linmaps:kerim}{Núcleo}}: escríbase $u(x) = 0$ como un sistema sobre las
\omterm{prop:b1:vspaces:coordinates}{coordenadas} (o los coeficientes) de $x$, resuélvase y
parametrícese --- el \omterm{def:b1:linmaps:kerim}{núcleo} sale con una \omterm{def:b1:vspaces:free}{base} incorporada
(\cref{met:b1:findim:computedim}). \emph{\omterm{def:b1:linmaps:kerim}{Imagen}}: es el
\omterm{def:b1:vspaces:span}{subespacio generado} por las imágenes de \emph{cualquier} familia
generadora de $E$ ---normalmente una \omterm{def:b1:vspaces:free}{base}, con lo que
$\operatorname{im} u = \operatorname{Vect}\bigl(u(e_1), \dots,
u(e_n)\bigr)$---; elimínense después las imágenes redundantes para
extraer una \omterm{def:b1:vspaces:free}{base}. \emph{Atajo}: calcúlese el más fácil de los dos y
obténgase gratis la dimensión del otro por el teorema del \omterm{def:b1:linmaps:rank}{rango}
(\cref{thm:b1:linmaps:ranknullity}); y cuando se conoce un
\omterm{def:b1:vspaces:subspace}{subespacio} candidato plausible para la \omterm{def:b1:linmaps:kerim}{imagen}, comparar
dimensiones eleva la inclusión fácil a igualdad
(\cref{thm:b1:findim:subspaces}). Los dos atajos se usan en
el~\cref{ex:b1:linmaps:delta} de más abajo.
\end{method}

\begin{proposition}\label{prop:b1:linmaps:kernel}
$u$ es \omterm{def:b1:logic:inj}{inyectiva} $\iff$ $\ker u = \{0\}$; y $u$ es \omterm{def:b1:logic:inj}{sobreyectiva}
$\iff$ $\operatorname{im} u = F$.
\end{proposition}

\begin{proof}
Como en los \omterm{def:b1:structures:group}{grupos} (\cref{prop:b1:structures:kernel}):
$u(x) = u(y) \iff u(x - y) = 0 \iff x - y \in \ker u$. El segundo
punto es la definición.
\end{proof}

\section{El teorema del rango}

\begin{definition}\label{def:b1:linmaps:rank}
El \emph{rango}\index{rango!de una aplicación lineal} de
$u \in \mathcal{L}(E, F)$ (con $E$ \omterm{def:b1:findim:def}{de dimensión finita}) es
$\operatorname{rk} u = \dim \operatorname{im} u$ --- también el rango
de la familia $\bigl(u(e_1), \dots, u(e_n)\bigr)$ para cualquier
base $(e_i)$ de $E$.
\end{definition}

\begin{theorem}[Teorema del rango]\label{thm:b1:linmaps:ranknullity}
Sea $E$ \omterm{def:b1:findim:def}{de dimensión finita} y $u \in \mathcal{L}(E, F)$.
Entonces\index{teorema del rango}
\[
\dim E = \dim \ker u + \operatorname{rk} u .
\]
Con más precisión: si $S$ es un \omterm{def:b1:vspaces:subspace}{subespacio} \omterm{def:b1:vspaces:sum}{suplementario} cualquiera de
$\ker u$ en $E$, entonces $u$ se restringe a un \emph{isomorfismo} de
$S$ sobre $\operatorname{im} u$.
\end{theorem}

\begin{proof}
Sea $S$ tal que $E = \ker u \oplus S$
(\cref{thm:b1:findim:subspaces}) y sea
$v \colon S \to \operatorname{im} u$ la restricción de $u$.

$v$ es \omterm{def:b1:logic:inj}{inyectiva}: $\ker v = S \cap \ker u = \{0\}$.

$v$ es \omterm{def:b1:logic:inj}{sobreyectiva}: cualquier $u(x)$ con $x = k + s$
($k \in \ker u$, $s \in S$) es igual a $u(s) = v(s)$.

Luego $v$ es un isomorfismo; y un isomorfismo lleva una \omterm{def:b1:vspaces:free}{base} a una
\omterm{def:b1:vspaces:free}{base} (\cref{prop:b1:linmaps:basis}), de modo que
$\dim S = \dim\operatorname{im} u$, y
$\dim E = \dim\ker u + \dim S$ concluye.
\end{proof}

\begin{example}[Construir una aplicación a medida]
\label{ex:b1:linmaps:tospec}
Constrúyase $u \in \mathcal{L}(\R^3)$ con
$\ker u = \operatorname{Vect}(1,1,1)$ y
$\operatorname{im} u = \{z = 0\}$. Primero, coherencia: el teorema
del \omterm{def:b1:linmaps:rank}{rango} exige $1 + 2 = 3$ --- cuadra, así que puede existir
solución. Elíjase una \omterm{def:b1:vspaces:free}{base} adaptada al \omterm{def:b1:linmaps:kerim}{núcleo}, digamos
$\bigl((1,1,1),\ e_1,\ e_2\bigr)$
(\cref{ex:b1:findim:completion}), y prescríbanse las imágenes
(\cref{prop:b1:linmaps:basis}):
\[
u(1,1,1) = 0, \qquad u(e_1) = e_1, \qquad u(e_2) = e_2 .
\]
Entonces $\ker u \supseteq \operatorname{Vect}(1,1,1)$ y
$\operatorname{im} u = \operatorname{Vect}(e_1, e_2) = \{z = 0\}$; y
el teorema del \omterm{def:b1:linmaps:rank}{rango} fuerza $\dim\ker u = 1$, de modo que el \omterm{def:b1:linmaps:kerim}{núcleo}
es exactamente la recta prescrita. Explícitamente, descomponiendo
$(x, y, z) = z(1,1,1) + (x - z)e_1 + (y - z)e_2$:
\[
u(x, y, z) = (x - z,\ y - z,\ 0).
\]
La receta se generaliza: existe una \omterm{def:b1:linmaps:def}{aplicación lineal} con \omterm{def:b1:linmaps:kerim}{núcleo}
prescrito $N$ e \omterm{def:b1:linmaps:kerim}{imagen} prescrita $I$ exactamente cuando
$\dim N + \dim I = \dim E$ --- la necesidad es el teorema del \omterm{def:b1:linmaps:rank}{rango} y
la suficiencia es esta construcción.
\end{example}

\begin{corollary}\label{cor:b1:linmaps:samedim}
Si $\dim E = \dim F$ (finita), entonces, para $u \in \mathcal{L}(E, F)$:
\[
u \text{ inyectiva} \iff u \text{ sobreyectiva} \iff u \text{
biyectiva}.
\]
En particular, esto vale para los endomorfismos en \omterm{def:b1:findim:def}{dimensión finita}.
(Falla en dimensión infinita: en $K[X]$, la \omterm{def:b1:derivative:def}{derivada} es \omterm{def:b1:logic:inj}{sobreyectiva}
pero no \omterm{def:b1:logic:inj}{inyectiva}, y $P \mapsto XP$ es \omterm{def:b1:logic:inj}{inyectiva} pero no
\omterm{def:b1:logic:inj}{sobreyectiva}.)
\end{corollary}

\begin{proof}
\omterm{def:b1:logic:inj}{Inyectiva} $\iff \dim\ker u = 0 \iff \operatorname{rk} u = \dim E =
\dim F \iff \operatorname{im} u = F$ (un \omterm{def:b1:vspaces:subspace}{subespacio} de dimensión
máxima lo es todo, \cref{thm:b1:findim:subspaces}) $\iff$
\omterm{def:b1:logic:inj}{sobreyectiva}.
\end{proof}

\begin{example}[La interpolación, estructuralmente]\label{ex:b1:linmaps:interpolation}
Fíjense $x_0, \dots, x_n$ distintos y sea
$u \colon \R_n[X] \to \R^{n+1}$,
$P \mapsto (P(x_0), \dots, P(x_n))$: \omterm{def:b1:linmaps:def}{lineal}. Su \omterm{def:b1:linmaps:kerim}{núcleo} es
$\{P : \deg P \leq n,\ n+1 \text{ raíces}\} = \{0\}$
(\cref{cor:b1:poly:nroots}). Dimensiones iguales $n + 1$: $u$ es un
isomorfismo --- la existencia \emph{y} la unicidad del interpolador
de Lagrange (\cref{thm:b1:poly:lagrange}) en una línea.

El mismo patrón de una línea trata datos que mezclan valores y
\omterm{def:b1:derivative:def}{derivadas}: $v \colon \R_3[X] \to \R^4$,
$P \mapsto \bigl(P(0), P'(0), P(1), P'(1)\bigr)$ es \omterm{def:b1:linmaps:def}{lineal}, y su
\omterm{def:b1:linmaps:kerim}{núcleo} consta de los \omterm{def:b1:poly:def}{polinomios} de grado $\leq 3$ con raíces dobles
en $0$ \emph{y} en $1$, es decir, divisibles por $X^2(X-1)^2$, de
grado $4$: solo $P = 0$. Otra vez dimensiones iguales: cada cuádrupla
de datos $(P(0), P'(0), P(1), P'(1))$ la realiza exactamente una
cúbica --- la interpolación de Hermite, concedida por un cálculo de
\omterm{def:b1:linmaps:kerim}{núcleo} antes de escribir ninguna fórmula (el problema del fin de
semana del~\cref{ch:b1:det} se encuentra con su determinante).
\end{example}

\begin{example}[El teorema del rango en acción: el operador de diferencias]
\label{ex:b1:linmaps:delta}
Sea $\Delta \colon \R_n[X] \to \R_n[X]$, $P \mapsto P(X+1) - P(X)$:
\omterm{def:b1:linmaps:def}{lineal}. \omterm{def:b1:linmaps:kerim}{Núcleo}: si $\Delta P = 0$, entonces
$P(0) = P(1) = P(2) = \dots$, luego $P - P(0)$ tiene infinitas raíces
y se anula (\cref{cor:b1:poly:nroots}): $\ker\Delta$ es la recta de
las constantes. Teorema del \omterm{def:b1:linmaps:rank}{rango}:
$\operatorname{rk}\Delta = (n + 1) - 1 = n$. Y como
$\deg \Delta P < \deg P$ para $P$ no constante (los términos de
cabeza se cancelan), $\operatorname{im}\Delta \subseteq \R_{n-1}[X]$,
que tiene dimensión exactamente $n$: la inclusión es una igualdad.
Conclusión, sin calcular ninguna \omterm{def:b1:logic:map}{imagen recíproca}: \emph{todo}
\omterm{def:b1:poly:def}{polinomio} $Q$ de grado $\leq n - 1$ es una diferencia
$Q = P(X+1) - P(X)$ --- la primitiva discreta existe. (Compárese con
el problema del fin de semana del~\cref{ch:b1:vspaces}, donde $\Delta$
se invirtió explícitamente en la \omterm{def:b1:vspaces:free}{base} binomial.)
\end{example}

\begin{example}[Contabilidad del rango a lo largo de una composición]
\label{ex:b1:linmaps:rankbookkeeping}
En $\R_2[X]$, compóngase la \omterm{def:b1:derivative:def}{derivada} $D(P) = P'$ (de \omterm{def:b1:linmaps:rank}{rango} $2$:
\omterm{def:b1:linmaps:kerim}{imagen} $\R_1[X]$, \omterm{def:b1:linmaps:kerim}{núcleo} las constantes) consigo misma. Entonces
$D \circ D = D^2$ lleva $P \mapsto P''$, con \omterm{def:b1:linmaps:kerim}{imagen} $\R_0[X]$:
\omterm{def:b1:linmaps:rank}{rango} $1$. Compárese con las cotas generales: la tosca da
$\operatorname{rk} D^2 \leq \min(2, 2) = 2$; y la fórmula exacta
del~\cref{exo:b1:linmaps:12} explica la pérdida con precisión,
\[
\operatorname{rk} D^2 = \operatorname{rk} D -
\dim\bigl(\ker D \cap \operatorname{im} D\bigr) = 2 - 1 = 1 ,
\]
ya que las constantes (\omterm{def:b1:linmaps:kerim}{núcleo} del $D$ exterior) están dentro de
$\R_1[X]$ (\omterm{def:b1:linmaps:kerim}{imagen} del $D$ \omterm{def:b1:topology:closure}{interior}) con dimensión $1$. El \omterm{def:b1:linmaps:rank}{rango} se
pierde exactamente donde el \omterm{def:b1:linmaps:kerim}{núcleo} exterior embosca a la \omterm{def:b1:linmaps:kerim}{imagen}
\omterm{def:b1:topology:closure}{interior} --- la frase que hay que recordar cuando los \omterm{def:b1:linmaps:rank}{rangos} de una
composición se portan mal.
\end{example}

\begin{example}[Núcleo e imagen del operador de Euler]
\label{ex:b1:linmaps:euler}
En $\R_n[X]$, sea $u(P) = X\,P'$ (\omterm{def:b1:linmaps:def}{lineal}: lo son derivar y
multiplicar por $X$). \emph{\omterm{def:b1:linmaps:kerim}{Núcleo}}: $XP' = 0$ fuerza $P' = 0$ (un
producto de \omterm{def:b1:poly:def}{polinomios} se anula solo si lo hace un factor), luego
$\ker u$ es la recta de las constantes. \emph{\omterm{def:b1:linmaps:kerim}{Imagen}}: sobre la
\omterm{def:b1:vspaces:free}{base} monomial,
\[
u(X^k) = k\,X^{k} \qquad (k = 0, 1, \dots, n),
\]
luego $\operatorname{im} u = \operatorname{Vect}(X, 2X^2, \dots,
nX^n) = \operatorname{Vect}(X, X^2, \dots, X^n)$: los \omterm{def:b1:poly:def}{polinomios} con
término constante nulo. Comprobación con el teorema del \omterm{def:b1:linmaps:rank}{rango}:
$\operatorname{rk} u = (n + 1) - 1 = n$, que es en efecto la
dimensión hallada. Dos observaciones que conviene guardar. Primera:
aquí $\operatorname{im} u \oplus \ker u = \R_n[X]$ --- pero eso es un
\emph{feliz accidente} de este operador, no un teorema: para el
$v(P) = P'$, de tipo desplazamiento, en $\R_1[X]$ se tiene
$\ker v = \operatorname{im} v = \R_0[X]$ y la suma no es directa.
Segunda: la relación $u(X^k) = kX^k$ dice que cada monomio
simplemente se reescala con $u$ --- una \omterm{def:b1:vspaces:free}{base} adaptada a la
\omterm{def:b1:logic:map}{aplicación}, el germen de la idea de valor propio que se desarrolla
en el volumen del segundo año.
\end{example}

\section{Proyecciones y simetrías}

\begin{definition}\label{def:b1:linmaps:projection}
Sea $E = F \oplus G$. La \emph{proyección sobre $F$ paralelamente a
$G$}\index{proyección} lleva $x = f + g$ (descomposición única) a
$p(x) = f$; la \emph{simetría}\index{simetría} asociada es
$s(x) = f - g$. Las dos son \omterm{def:b1:linmaps:def}{lineales}, y $s = 2p - \mathrm{id}$.
\end{definition}

\begin{theorem}[Algebraic characterization]\label{thm:b1:linmaps:projchar}
\begin{enumerate}
  \item Un endomorfismo $p$ es una \omterm{def:b1:linmaps:projection}{proyección} (sobre cierto $F$
        paralelamente a cierto $G$) si y solo si $p \circ p = p$; y
        entonces $F = \operatorname{im} p = \ker(p - \mathrm{id})$ y
        $G = \ker p$.
  \item Un endomorfismo $s$ es una \omterm{def:b1:linmaps:projection}{simetría} si y solo si
        $s \circ s = \mathrm{id}$; y entonces
        $E = \ker(s - \mathrm{id}) \oplus \ker(s + \mathrm{id})$.
\end{enumerate}
\end{theorem}

\begin{proof}
(1) Una \omterm{def:b1:linmaps:projection}{proyección} cumple $p(f + g) = f$ y $p(f) = f$: $p^2 = p$.
Recíprocamente, sea $p^2 = p$; póngase $F = \operatorname{im} p$,
$G = \ker p$. Todo $x$ se escribe $x = p(x) + (x - p(x))$ con
$p(x) \in F$ y $p\bigl(x - p(x)\bigr) = p(x) - p^2(x) = 0$:
$E = F + G$. Si $y \in F \cap G$: $y = p(z)$ y $p(y) = 0$, luego
$y = p(z) = p^2(z) = p(y) = 0$: \omterm{def:b1:vspaces:sum}{suma directa}, y $p$ es la
\omterm{def:b1:linmaps:projection}{proyección} sobre $F$ paralelamente a $G$. Por último, en $F$:
$y = p(z)$ da $p(y) = y$, luego $F \subseteq \ker(p - \mathrm{id})$;
y, recíprocamente, $p(y) = y$ mete a $y$ en la \omterm{def:b1:linmaps:kerim}{imagen}.

(2) La correspondencia $s = 2p - \mathrm{id}$,
$p = \frac{s + \mathrm{id}}2$ es una \omterm{def:b1:logic:inj}{biyección} entre endomorfismos, y
mediante ella
\[
s^2 = 4p^2 - 4p + \mathrm{id} = \mathrm{id}
\iff 4p^2 = 4p \iff p^2 = p :
\]
las \omterm{def:b1:linmaps:projection}{simetrías} se corresponden exactamente con las proyecciones.
Traduciendo los \omterm{def:b1:vspaces:subspace}{subespacios}: $s(x) = x \iff p(x) = x$, luego
$\ker(s - \mathrm{id}) = \operatorname{im} p = F$; y
$s(x) = -x \iff 2p(x) = 0 \iff x \in \ker p = G$, luego
$\ker(s + \mathrm{id}) = G$. La \omterm{def:b1:vspaces:sum}{suma directa} $E = F \oplus G$ del
punto (1) se convierte en la descomposición anunciada en los vectores
fijos y los vectores invertidos por $s$.
\end{proof}

\begin{example}[Una proyección y su simetría, explícitamente]
\label{ex:b1:linmaps:projexplicit}
En $\R^2$, proyéctese sobre $F = \operatorname{Vect}(1,1)$
paralelamente a $G = \operatorname{Vect}(0,1)$. Descompóngase
$(x, y) = a(1,1) + b(0,1)$: la primera coordenada da $a = x$ y la
segunda, $b = y - x$. Por tanto,
\[
p(x, y) = (x, x),
\qquad
s(x, y) = 2p(x,y) - (x,y) = (x,\ 2x - y).
\]
Compruébese el álgebra: $p(p(x,y)) = p(x,x) = (x,x)$ y
$s(s(x,y)) = s(x, 2x - y) = (x, 2x - (2x - y)) = (x, y)$.
Geométricamente, $s$ es la «reflexión oblicua» respecto de la recta
$y = x$ en la dirección vertical: deja $F$ fijo punto a punto e
invierte $G$. Si hubiéramos proyectado sobre el mismo $F$ pero
paralelamente a $G' = \operatorname{Vect}(1,-1)$, la fórmula
cambiaría a $p'(x,y) = \bigl(\frac{x+y}2, \frac{x+y}2\bigr)$: una
\omterm{def:b1:linmaps:projection}{proyección} queda determinada por su \omterm{def:b1:linmaps:kerim}{imagen} \emph{y} su \omterm{def:b1:linmaps:kerim}{núcleo},
nunca por la \omterm{def:b1:linmaps:kerim}{imagen} sola.
\end{example}

\begin{omfigure}
\begin{tikzpicture}[scale=1.15]
  \draw[->, thin] (-0.4,0) -- (3.4,0) node[below] {$x$};
  \draw[->, thin] (0,-0.4) -- (0,3.8) node[left] {$y$};
  \draw[omDef, very thick] (-0.3,-0.3) -- (3.1,3.1)
    node[right, pos=0.98, font=\small] {$F:\ y = x$};
  \draw[omProp, thick, ->] (2.6,0.2) -- (2.6,0.9)
    node[right, font=\small] {$G$};
  \fill (2,0.5) circle (0.05) node[right, font=\small] {$M$};
  \fill (2,2) circle (0.05) node[left, font=\small] {$p(M)$};
  \fill (2,3.5) circle (0.05) node[left, font=\small] {$s(M)$};
  \draw[omThm, dashed] (2,0.5) -- (2,3.5);
\end{tikzpicture}

{\small La \omterm{def:b1:linmaps:projection}{proyección} sobre $F = \operatorname{Vect}(1,1)$
paralelamente a $G = \operatorname{Vect}(0,1)$ y su \omterm{def:b1:linmaps:projection}{simetría}, sobre
el punto $M = (2,\ 0.5)$: deslizándose verticalmente, $M$ llega a $F$
en $p(M) = (2,2)$ y aterriza en $s(M) = 2p(M) - M = (2,\ 3.5)$, tan
por encima de $F$ (medido a lo largo de $G$) como $M$ estaba por
debajo.}
\end{omfigure}

\section{Formas lineales e hiperplanos}

\begin{definition}\label{def:b1:linmaps:forms}
Una \emph{forma lineal}\index{forma lineal} en $E$ es una
\omterm{def:b1:linmaps:def}{aplicación lineal} $\varphi \colon E \to K$. Un
\emph{hiperplano}\index{hiperplano} de $E$ ($\dim E = n$) es un
\omterm{def:b1:vspaces:subspace}{subespacio} de dimensión $n - 1$.
\end{definition}

\begin{example}[Una forma de evaluación y su hiperplano]
\label{ex:b1:linmaps:evaluationform}
En $\R_2[X]$, la evaluación $\varphi(P) = P(2)$ es una \omterm{def:b1:linmaps:forms}{forma lineal}
no nula ($\varphi(1) = 1$). Su \omterm{def:b1:linmaps:kerim}{núcleo} es el \omterm{def:b1:linmaps:forms}{hiperplano} de los
\omterm{def:b1:poly:def}{polinomios} que se anulan en $2$, es decir (teorema del factor,
\cref{thm:b1:poly:factor}), los múltiplos de $X - 2$ dentro de
$\R_2[X]$:
\[
\ker\varphi = \operatorname{Vect}\bigl(X - 2,\ X(X - 2)\bigr),
\qquad \dim = 2 .
\]
En \omterm{prop:b1:vspaces:coordinates}{coordenadas} sobre $(1, X, X^2)$,
$\varphi(a + bX + cX^2) = a + 2b + 4c$: toda \omterm{def:b1:linmaps:forms}{forma lineal} en un
espacio \omterm{def:b1:findim:def}{de dimensión finita} es, una vez fijada una \omterm{def:b1:vspaces:free}{base}, una
expresión \omterm{def:b1:linmaps:def}{lineal} fija en las \omterm{prop:b1:vspaces:coordinates}{coordenadas} --- las formas son
«vectores fila», como el~\cref{ch:b1:matrices} volverá literal, y la
fila de coeficientes de aquí, $(1, 2, 4)$, es una fila de
Vandermonde: las formas de evaluación son la puerta por la que entra
en el álgebra lineal la teoría de interpolación del problema del fin
de semana del~\cref{ch:b1:det}.
\end{example}

\begin{theorem}\label{thm:b1:linmaps:hyperplanes}
Los \omterm{def:b1:linmaps:forms}{hiperplanos} de $E$ son exactamente los \omterm{def:b1:linmaps:kerim}{núcleos} de las
\omterm{def:b1:linmaps:forms}{formas lineales} no nulas. Y dos formas no nulas tienen el mismo
\omterm{def:b1:linmaps:kerim}{núcleo} si y solo si son proporcionales.
\end{theorem}

\begin{proof}
Si $\varphi \neq 0$: $\operatorname{rk}\varphi = 1$ (la \omterm{def:b1:linmaps:kerim}{imagen} es un
\omterm{def:b1:vspaces:subspace}{subespacio} no nulo de $K$), luego $\dim\ker\varphi = n - 1$: un
\omterm{def:b1:linmaps:forms}{hiperplano}. Recíprocamente, sea $H$ un \omterm{def:b1:linmaps:forms}{hiperplano} y
$(e_1, \dots, e_{n-1})$ una \omterm{def:b1:vspaces:free}{base} de $H$ completada por $e_n$: la
forma «última coordenada» tiene \omterm{def:b1:linmaps:kerim}{núcleo} $H$.

Dos formas proporcionales comparten su \omterm{def:b1:linmaps:kerim}{núcleo}. Recíprocamente,
supóngase $\ker\varphi = \ker\psi = H$ y tómese $a \notin H$: todo
$x$ se escribe $x = h + \lambda a$ (pues $E = H \oplus Ka$), y
\[
\varphi(x) = \lambda \varphi(a), \qquad \psi(x) = \lambda\psi(a):
\]
luego $\varphi = \frac{\varphi(a)}{\psi(a)}\,\psi$.
\end{proof}

\begin{example}\label{ex:b1:linmaps:hyperplane}
En $K^n$, un \omterm{def:b1:linmaps:forms}{hiperplano} es un \omterm{def:b1:logic:sets}{conjunto} de soluciones
$\{a_1 x_1 + \dots + a_n x_n = 0\}$ con los $a_i$ no todos nulos ---
la conocida ecuación de un plano que pasa por el origen en $\R^3$. En
los espacios de funciones, las formas de evaluación $P \mapsto P(1)$
o $f \mapsto \int_0^1 f$ definen \omterm{def:b1:linmaps:forms}{hiperplanos} de $\R_n[X]$ y de
$C(\intcc{0}{1})$ (cf.\ \cref{exo:b1:findim:6}).
\end{example}

\begin{example}[Un hiperplano, trabajado de tres maneras]
\label{ex:b1:linmaps:hyperplanerun}
Tómense $\varphi(x, y, z) = x - 2y + 3z$ en $\R^3$ y
$H = \ker\varphi$. \emph{\omterm{def:b1:vspaces:free}{Base}}: resuélvase $x = 2y - 3z$:
\[
(2y - 3z,\ y,\ z) = y\,(2, 1, 0) + z\,(-3, 0, 1),
\]
dos vectores \omterm{def:b1:vspaces:free}{libres}: $\dim H = 2$, un \omterm{def:b1:linmaps:forms}{hiperplano}, como predice
el~\cref{thm:b1:linmaps:hyperplanes} a partir de $\varphi \neq 0$.
\emph{Recta suplementaria}: cualquier vector fuera de $H$ genera
una, por ejemplo $a = (1, 0, 0)$ ($\varphi(a) = 1 \neq 0$); y la
descomposición de un $v$ arbitrario es explícita:
\[
v = \underbrace{\bigl(v - \varphi(v)\,a\bigr)}_{\in\,H}
+ \underbrace{\varphi(v)\,a}_{\in\,\operatorname{Vect}(a)},
\]
ya que $\varphi\bigl(v - \varphi(v)a\bigr) = \varphi(v) -
\varphi(v)\varphi(a) = 0$. \emph{Proporcionalidad}: si
$\psi(x,y,z) = -2x + 4y - 6z$, entonces $\psi = -2\varphi$ y las dos
tienen \omterm{def:b1:linmaps:kerim}{núcleo} $H$; recíprocamente, toda forma que se anule en $H$
es múltiplo de $\varphi$ (\cref{exo:b1:linmaps:8}) --- la ecuación de
un \omterm{def:b1:linmaps:forms}{hiperplano} es única salvo escala, hecho que se usa
constantemente para los planos en geometría.
\end{example}

\begin{remark}[Errores frecuentes]
\label{rem:b1:linmaps:pitfalls}
\emph{El \omterm{def:b1:linmaps:kerim}{núcleo} y la \omterm{def:b1:linmaps:kerim}{imagen} viven en espacios distintos}:
$\ker u \subseteq E$, $\operatorname{im} u \subseteq F$; la suma
$\ker u + \operatorname{im} u$ solo tiene sentido para
endomorfismos y, aun así, no tiene por qué ser directa
($u(x, y) = (y, 0)$ tiene $\ker u = \operatorname{im} u$; el
\cref{exo:b1:linmaps:7} caracteriza cuándo sí lo es).
\emph{$u^2 = 0$ no significa $u = 0$}: ese mismo $u(x,y) = (y, 0)$
tiene cuadrado nulo sin anularse --- lo que $u^2 = 0$ dice de verdad
es $\operatorname{im} u \subseteq \ker u$
(\cref{exo:b1:linmaps:5}).
\emph{\omterm{def:b1:logic:inj}{Inyectiva} $\iff$ \omterm{def:b1:logic:inj}{sobreyectiva} exige dimensiones finitas
iguales}: en $K[X]$, la \omterm{def:b1:derivative:def}{derivada} es \omterm{def:b1:logic:inj}{sobreyectiva} y no \omterm{def:b1:logic:inj}{inyectiva}, y
$P \mapsto XP$ es \omterm{def:b1:logic:inj}{inyectiva} y no \omterm{def:b1:logic:inj}{sobreyectiva}
(\cref{cor:b1:linmaps:samedim}); y entre espacios de dimensiones
\emph{distintas}, una de las dos implicaciones es sencillamente
imposible ($\operatorname{rk} u \leq \min(\dim E, \dim F)$).
\emph{Prescribir imágenes funciona sobre una \omterm{def:b1:vspaces:free}{base}, no sobre una
familia cualquiera}: exigir $u(1, 0) = a$, $u(0, 1) = b$,
$u(1, 1) = c$ sobredetermina $u$ salvo que $c = a + b$; una
\omterm{def:b1:linmaps:def}{aplicación lineal} es \omterm{def:b1:vspaces:free}{libre} en una \omterm{def:b1:vspaces:free}{base} y esclava en todo lo demás.
\emph{El \omterm{def:b1:linmaps:rank}{rango} no se conserva al componer}: solo puede bajar,
$\operatorname{rk}(vu) \leq \min(\operatorname{rk} u,
\operatorname{rk} v)$ (\cref{exo:b1:linmaps:4}), con la pérdida
exacta medida en el~\cref{exo:b1:linmaps:12}.
\end{remark}

\begin{remark}[Adónde van estas aplicaciones]
\label{rem:b1:linmaps:whereused}
Las \omterm{def:b1:linmaps:def}{aplicaciones lineales} están a punto de convertirse en
\emph{matrices}: una vez fijadas las \omterm{def:b1:vspaces:free}{bases},
el~\cref{ch:b1:matrices} codifica cada $u \in \mathcal{L}(E, F)$
mediante una tabla rectangular, y la composición pasa a ser el
producto de matrices --- y el teorema del \omterm{def:b1:linmaps:rank}{rango} dirige después la
teoría de los sistemas \omterm{def:b1:linmaps:def}{lineales} en el~\cref{ch:b1:det}. Las
proyecciones vuelven en el~\cref{ch:b1:euclid} en su caso
particular más útil, la \omterm{def:b1:linmaps:projection}{proyección} \emph{ortogonal}, donde el
\omterm{def:b1:linmaps:kerim}{núcleo} se elige perpendicular a la \omterm{def:b1:linmaps:kerim}{imagen}. El problema del fin de
semana de más abajo lleva el álgebra de proyectores tan lejos como
alcanzan las herramientas de primer año, hasta el lema de Fitting; y
el volumen del segundo año va más allá, con la traza y con la teoría
de valores propios, para la que los proyectores sobre \omterm{def:b1:vspaces:subspace}{subespacios}
estables son las piezas básicas.
\end{remark}

\begin{remark}[Perspectivas dentro del libro 3: el teorema del rango tres veces más]
\label{rem:b1:linmaps:perspectives}
La ley de conservación $\dim E = \dim\ker u + \operatorname{rk} u$ se
releerá tres veces antes de que acabe el volumen. En
el~\cref{ch:b1:det} se convierte en la forma de los \omterm{def:b1:logic:sets}{conjuntos} de
soluciones: un sistema compatible con $p$ incógnitas y \omterm{def:b1:linmaps:rank}{rango} $r$
tiene un \omterm{def:b1:logic:sets}{conjunto} de soluciones de dimensión $p - r$ --- la dimensión
del \omterm{def:b1:linmaps:kerim}{núcleo} disfrazada. En el~\cref{ch:b1:euclid} se parte de forma
ortogonal, $\dim F + \dim F^\perp = \dim E$, y mueve todo cálculo de
distancias. Y en el problema del fin de semana
del~\cref{ch:b1:multivar} es el contable de los mínimos cuadrados:
$n$ observaciones, $2$ parámetros ajustados, $n - 2$ dimensiones de
residuo, y la identidad pitagórica
$\norm b^2 = \norm p^2 + \norm{b - p}^2$ es la sombra euclídea del
teorema del \omterm{def:b1:linmaps:rank}{rango}. Un teorema, cuatro disfraces.
\end{remark}

\section{Ejercicios}

\begin{exercise}[$\star$]\label{exo:b1:linmaps:1}
¿Qué aplicaciones son \omterm{def:b1:linmaps:def}{lineales}?
\begin{enumerate}
  \item $\R^2 \to \R^2$, $(x, y) \mapsto (x + y, x - 2y)$;
  \item $\R^2 \to \R$, $(x, y) \mapsto xy$;
  \item $\R[X] \to \R[X]$, $P \mapsto P' + XP$;
  \item $\mathcal{F}(\R,\R) \to \R$, $f \mapsto f(3)$.
\end{enumerate}
\end{exercise}

\begin{exercise}[$\star$]\label{exo:b1:linmaps:2}
Sea $u \colon \R^3 \to \R^3$,
$(x,y,z) \mapsto (x + y - z,\; 2x + y + z,\; 3x + 2y)$. Determínense
$\ker u$ (\omterm{def:b1:vspaces:free}{base} y dimensión), $\operatorname{rk} u$ y una \omterm{def:b1:vspaces:free}{base} de
$\operatorname{im} u$. ¿Es $u$ \omterm{def:b1:logic:inj}{inyectiva}? ¿\omterm{def:b1:logic:inj}{sobreyectiva}?
\end{exercise}

\begin{exercise}[$\star$]\label{exo:b1:linmaps:3}
Sea $u \colon \R_n[X] \to \R_n[X]$, $P \mapsto P - P'$. Demuéstrese
que $u$ es un isomorfismo: una vez por $\ker u$ y otra exhibiendo la
inversa \emph{(considérese $P + P' + P'' + \dots$)}.
\end{exercise}

\begin{exercise}[$\star$]\label{exo:b1:linmaps:4}
Sean $u \in \mathcal{L}(E, F)$ y $v \in \mathcal{L}(F, G)$, con los
espacios \omterm{def:b1:findim:def}{de dimensión finita}. Demuéstrese:
\[
\operatorname{rk}(v \circ u) \leq
\min\bigl(\operatorname{rk} u,\ \operatorname{rk} v\bigr).
\]
\end{exercise}

\begin{exercise}[$\star\star$]\label{exo:b1:linmaps:5}
Sea $u$ un endomorfismo de $E$ (\omterm{def:b1:findim:def}{de dimensión finita}) con $u^2 = 0$.
Demuéstrese que $\operatorname{im} u \subseteq \ker u$ y, por tanto,
que $\operatorname{rk} u \leq \frac{\dim E}{2}$. Para $E = \R^2$,
dese un ejemplo con igualdad.
\end{exercise}

\begin{exercise}[$\star\star$]\label{exo:b1:linmaps:6}
Sean $p, q$ proyecciones de $E$ con $p \circ q = q \circ p$.
Demuéstrese que $p \circ q$ es una \omterm{def:b1:linmaps:projection}{proyección}, con
\[
\operatorname{im}(pq) = \operatorname{im} p \cap \operatorname{im}
q ,
\qquad
\ker (pq) = \ker p + \ker q .
\]
\end{exercise}

\begin{exercise}[$\star\star$]\label{exo:b1:linmaps:7}
Sea $u \in \mathcal{L}(E)$, con $E$ \omterm{def:b1:findim:def}{de dimensión finita}.
Demuéstrese la equivalencia de:
\begin{enumerate}
  \item $E = \ker u \oplus \operatorname{im} u$;
  \item $\ker u = \ker u^2$;
  \item $\operatorname{im} u = \operatorname{im} u^2$.
\end{enumerate}
\end{exercise}

\begin{exercise}[$\star\star$]\label{exo:b1:linmaps:8}
Sean $\varphi, \psi$ \omterm{def:b1:linmaps:forms}{formas lineales} en $E$ con
$\ker\varphi \subseteq \ker\psi$. Demuéstrese que
$\psi = \lambda\varphi$ para algún $\lambda \in K$ (incluidos los
casos degenerados).
\end{exercise}

\begin{exercise}[$\star\star\star$]\label{exo:b1:linmaps:9}
Sea $u \in \mathcal{L}(E)$ con $\dim E = n$, y supóngase $u^n = 0$
pero $u^{n-1} \neq 0$ (un endomorfismo \emph{nilpotente máximo}).
Tómese $x$ con $u^{n-1}(x) \neq 0$; demuéstrese que
$\bigl(x, u(x), \dots, u^{n-1}(x)\bigr)$ es una \omterm{def:b1:vspaces:free}{base} de $E$.
\emph{(Aplíquense potencias de $u$ a una combinación nula, empezando
por $u^{n-1}$.)}
\end{exercise}

\begin{exercise}[$\star\star\star$]\label{exo:b1:linmaps:10}
Sea $f \in \mathcal{L}(\R^n)$ con $f \circ f = -\mathrm{id}$.
\begin{enumerate}
  \item Demuéstrese que $f$ es un isomorfismo y que ningún
        $x \neq 0$ cumple $f(x) = \lambda x$ con $\lambda \in \R$.
  \item Demuéstrese que $n$ es par. \emph{Indicación: tómese
        $x_1 \neq 0$; véase que $\operatorname{Vect}(x_1, f(x_1))$
        es un plano estable por $f$; elíjase $x_2$ fuera de él e
        itérese, demostrando que
        $\bigl(x_1, f(x_1), x_2, f(x_2), \dots\bigr)$ sigue \omterm{def:b1:vspaces:free}{libre}.}
\end{enumerate}
\end{exercise}

\begin{exercise}[$\star\star$]\label{exo:b1:linmaps:11}
Sean $u, v \in \mathcal{L}(E, F)$, con los espacios \omterm{def:b1:findim:def}{de
dimensión finita}. Demuéstrese la cota por los dos lados
\[
\abs{\operatorname{rk} u - \operatorname{rk} v}
\;\leq\; \operatorname{rk}(u + v)
\;\leq\; \operatorname{rk} u + \operatorname{rk} v .
\]
\emph{(Para la \omterm{def:b1:reals:bounds}{cota superior}, compárese $\operatorname{im}(u+v)$ con
$\operatorname{im} u + \operatorname{im} v$; y para la inferior,
aplíquese hábilmente la \omterm{def:b1:reals:bounds}{cota superior}.)}
\end{exercise}

\begin{exercise}[$\star\star\star$]\label{exo:b1:linmaps:12}
(Desigualdad de Frobenius) Sean $u \in \mathcal{L}(E, F)$,
$w \in \mathcal{L}(F, G)$ y $v \in \mathcal{L}(G, H)$, con todos los
espacios \omterm{def:b1:findim:def}{de dimensión finita}. Demuéstrese la fórmula exacta
\[
\operatorname{rk}(v \circ w) = \operatorname{rk} w -
\dim\bigl(\ker v \cap \operatorname{im} w\bigr),
\]
y dedúzcase la desigualdad de Frobenius
\[
\operatorname{rk}(v \circ w) + \operatorname{rk}(w \circ u)
\;\leq\; \operatorname{rk} w + \operatorname{rk}(v \circ w \circ
u) .
\]
Compruébese que el caso $w = \mathrm{id}_F$ es la desigualdad de
Sylvester, demostrada en forma matricial en el~\cref{exo:b1:matrices:10}.
\end{exercise}

\section{Problema: cálculo de proyectores y lema de Fitting}

\begin{problem}\label{pb:b1:linmaps:1}
Las proyecciones son los endomorfismos que producen las \omterm{def:b1:vspaces:sum}{sumas
directas}, y recíprocamente: toda identidad
$E = F_1 \oplus \dots \oplus F_k$ es en secreto una familia de
proyectores que suman la identidad. Este problema desarrolla ese
diccionario ---el álgebra de un proyector, de dos, de $k$--- y aplica
después las mismas ideas de estabilización a un endomorfismo
arbitrario, demostrando el \emph{lema de Fitting}: todo endomorfismo
de un espacio \omterm{def:b1:findim:def}{de dimensión finita} se parte en una parte nilpotente y
una parte invertible. En todo el problema, $E$ es un $K$-espacio
vectorial de dimensión $n$, y \emph{proyector} significa
$p \in \mathcal{L}(E)$ con $p^2 = p$
(\cref{thm:b1:linmaps:projchar}).
\index{proyección}\index{lema de Fitting}\index{nilpotente}

\medskip
\noindent\textbf{Parte I --- El álgebra en torno a un proyector.}
Sea $p$ un proyector, $p \neq 0$, $p \neq \mathrm{id}$.
\begin{enumerate}
  \item Véase que $\mathrm{id} - p$ es un proyector e identifíquense
        $\operatorname{im}(\mathrm{id} - p)$ y
        $\ker(\mathrm{id} - p)$.
  \item Calcúlese $(\lambda\,\mathrm{id} + \mu\,p)^2$ y
        determínense todas las parejas $(\lambda, \mu) \in K^2$ para
        las que $\lambda\,\mathrm{id} + \mu\,p$ es un proyector.
  \item Véase que el plano $\operatorname{Vect}(\mathrm{id}, p)$ de
        $\mathcal{L}(E)$ es estable por composición y que, para todo
        \omterm{def:b1:poly:def}{polinomio} $Q \in K[X]$,
        \[
        Q(p) = Q(0)\,\mathrm{id} + \bigl(Q(1) -
        Q(0)\bigr)\,p .
        \]
  \item Determínese para qué $(\lambda, \mu)$ la \omterm{def:b1:logic:map}{aplicación}
        $\lambda\,\mathrm{id} + \mu\,p$ es invertible, y dese su
        inversa en la forma $\alpha\,\mathrm{id} + \beta\,p$.
        Interprétese la respuesta mediante la acción de
        $\lambda\,\mathrm{id} + \mu\,p$ sobre $\operatorname{im} p$ y
        sobre $\ker p$.
  \item Sea $p'$ otro proyector con la \emph{misma \omterm{def:b1:linmaps:kerim}{imagen}}
        $\operatorname{im} p' = \operatorname{im} p$. Véase que
        $p\,p' = p'$ y $p'\,p = p$. ¿Qué dicen estas identidades
        sobre componer proyecciones sobre el mismo \omterm{def:b1:vspaces:subspace}{subespacio}
        paralelamente a \omterm{def:b1:linmaps:kerim}{núcleos} distintos?
\end{enumerate}

\medskip
\noindent\textbf{Parte II --- Dos proyectores.} Sean $p, q$
proyectores de $E$; supóngase que la característica no es $2$ (cierto
para $K = \R, \C$).
\begin{enumerate}[resume]
  \item Supóngase que $p + q$ es un proyector. Desarrollando
        $(p + q)^2$, véase que $pq + qp = 0$; componiendo con $p$ por
        la izquierda y después por la derecha, dedúzcase $pq = qp$ y
        conclúyase $pq = qp = 0$.
  \item Recíprocamente, supóngase $pq = qp = 0$. Véase que $p + q$ es
        un proyector, con
        \[
        \operatorname{im}(p + q) = \operatorname{im} p \oplus
        \operatorname{im} q,
        \qquad
        \ker(p + q) = \ker p \cap \ker q .
        \]
  \item Véase que $p - q$ es un proyector si y solo si
        $pq = qp = q$. \emph{(Aplíquense las preguntas 6--7 a
        $\mathrm{id} - p$ y a $q$.)}
  \item Véase el significado geométrico de $pq = qp = q$: se cumple
        si y solo si $\operatorname{im} q \subseteq
        \operatorname{im} p$ y $\ker p \subseteq \ker q$. (Entonces
        se escribe $q \leq p$: «$q$ proyecta sobre menos,
        paralelamente a más».)
  \item Supóngase ahora que $p, q$ conmutan. Recuérdese
        del~\cref{exo:b1:linmaps:6} que $pq$ es el proyector sobre
        $\operatorname{im} p \cap \operatorname{im} q$
        paralelamente a $\ker p + \ker q$. Véase que
        $r = p + q - pq$ es un proyector con
        \[
        \operatorname{im} r = \operatorname{im} p +
        \operatorname{im} q,
        \qquad
        \ker r = \ker p \cap \ker q .
        \]
        \emph{(Considérese $\mathrm{id} - r = (\mathrm{id} -
        p)(\mathrm{id} - q)$.)}
\end{enumerate}

\medskip
\noindent\textbf{Parte III --- Descomposiciones de la identidad.}
\begin{enumerate}[resume]
  \item Sea $E = F_1 \oplus \dots \oplus F_k$ y, para
        $x = x_1 + \dots + x_k$ (descomposición única,
        $x_i \in F_i$), póngase $p_i(x) = x_i$. Véase que cada $p_i$
        es un proyector, que $p_i p_j = 0$ para $i \neq j$ y que
        $p_1 + \dots + p_k = \mathrm{id}$; identifíquense
        $\operatorname{im} p_i$ y $\ker p_i$.
  \item Recíprocamente, sean $p_1, \dots, p_k \in \mathcal{L}(E)$
        con $p_1 + \dots + p_k = \mathrm{id}$ y $p_i p_j = 0$ para
        todos $i \neq j$. Véase que cada $p_i$ es un proyector y que
        $E = \operatorname{im} p_1 \oplus \dots \oplus
        \operatorname{im} p_k$.
  \item Dos proyectores con $p + q = \mathrm{id}$: véase que
        $pq = qp = 0$ se cumple automáticamente.
  \item Tres proyectores con $p + q + r = \mathrm{id}$: véase que
        $p + q$ es un proyector y dedúzcase de la pregunta 6 que se
        anulan \emph{todos} los productos dos a dos --- luego
        $E = \operatorname{im} p \oplus \operatorname{im} q \oplus
        \operatorname{im} r$, sin ninguna hipótesis sobre los
        productos.
  \item Para $k$ proyectores con $p_1 + \dots + p_k = \mathrm{id}$:
        véase primero que, para \omterm{def:b1:vspaces:subspace}{subespacios} cualesquiera,
        $\dim (F_1 + \dots + F_k) \leq \dim F_1 + \dots + \dim F_k$,
        con igualdad si y solo si la suma es directa; véase después
        que $E = \operatorname{im} p_1 + \dots + \operatorname{im}
        p_k$, y demuéstrese que \emph{si} además
        $\sum_i \operatorname{rk} p_i \leq n$, la suma es directa y
        $p_i p_j = 0$ para $i \neq j$.
\end{enumerate}

\medskip
\noindent\textbf{Parte IV --- \omterm{def:b1:linmaps:kerim}{Núcleos} iterados: el lema de Fitting.}
Sea $u \in \mathcal{L}(E)$, $\dim E = n$.
\begin{enumerate}[resume]
  \item Véanse las dos cadenas, válidas para todo $k \geq 0$:
        \[
        \ker u^k \subseteq \ker u^{k+1},
        \qquad
        \operatorname{im} u^{k+1} \subseteq \operatorname{im}
        u^k .
        \]
  \item Véase que si $\ker u^{r} = \ker u^{r+1}$ para algún $r$,
        entonces $\ker u^{k} = \ker u^{r}$ para todo $k \geq r$;
        enúnciese y demuéstrese la estabilización análoga para las
        imágenes.
  \item Dedúzcase que hay un entero mínimo $r$ con
        $\ker u^{r} = \ker u^{r+1}$, que $r \leq n$ y que las
        imágenes se estabilizan en ese mismo $r$.
  \item (Lema de Fitting) Demuéstrese que
        \[
        E \;=\; \ker u^{r} \,\oplus\, \operatorname{im} u^{r} .
        \]
  \item Véase que los dos \omterm{def:b1:vspaces:subspace}{subespacios} son estables por $u$, que la
        restricción de $u$ a $\ker u^{r}$ es nilpotente y que la
        restricción de $u$ a $\operatorname{im} u^{r}$ es un
        isomorfismo de $\operatorname{im} u^{r}$: todo endomorfismo
        es, sobre una \omterm{def:b1:vspaces:sum}{suma directa} canónica, «nilpotente más
        invertible».
  \item Sea $\pi$ el proyector sobre $\ker u^{r}$ paralelamente a
        $\operatorname{im} u^{r}$. Véase que
        $\pi \circ u = u \circ \pi$.
\end{enumerate}

\medskip
\noindent\textbf{Parte V --- Un caso trabajado, y síntesis.}
\begin{enumerate}[resume]
  \item Sea $u(x, y, z) = (y, 0, z)$ en $\R^3$. Calcúlense $u^2$ y
        $u^3$, determínense el índice de estabilización $r$, los
        \omterm{def:b1:vspaces:subspace}{subespacios} $\ker u^{r}$ y $\operatorname{im} u^{r}$ y el
        proyector de Fitting $\pi$, y compruébese en las fórmulas que
        $\pi u = u\pi$ y que $u$ es nilpotente en un factor y
        \omterm{def:b1:logic:inj}{biyectiva} en el otro.
  \item Véanse las equivalencias: $u$ nilpotente $\iff$
        $\ker u^{r} = E$ $\iff$ $\pi = \mathrm{id}$; y dedúzcase que
        un endomorfismo nilpotente de un espacio de dimensión $n$
        siempre cumple $u^{n} = 0$ (el índice de nilpotencia nunca
        supera la dimensión).
  \item (Unicidad) Supóngase $E = A \oplus B$ con $A, B$ estables
        por $u$, la restricción $u|_A$ nilpotente y $u|_B$
        \omterm{def:b1:logic:inj}{biyectiva}. Demuéstrese $A = \ker u^{r}$ y
        $B = \operatorname{im} u^{r}$: la descomposición de Fitting
        es única.
  \item Síntesis, en cuatro frases: qué diccionario establece la
        parte III entre las \omterm{def:b1:vspaces:sum}{sumas directas} y las familias de
        proyectores; por qué la pregunta 14 no necesitó hipótesis
        sobre los productos mientras que la 15 necesitó una
        hipótesis de \omterm{def:b1:linmaps:rank}{rango} (y qué herramienta del segundo año, la
        traza, la elimina); en qué sentido el lema de Fitting es la
        versión estabilizada del~\cref{exo:b1:linmaps:7}; y en qué se
        convierten los dos factores de Fitting en la teoría de
        valores propios del volumen del segundo año. Nómbrese el
        teorema demostrado en la parte IV.
\end{enumerate}
\end{problem}
